% Run GCT_KRON_20261120; GPT-6 Astra Pro; effort not exposed.
\section{A second-row-sensitive vanishing theorem for Chow equations}

\subsection*{A-020. A shape-dependent improvement of the degree gate}
Let $V$ be a finite-dimensional complex vector space, and let
$\operatorname{Ch}_b(V)\subseteq\operatorname{Sym}^bV$ be the
Zariski closure of products of $b$ vectors, for $b\geq1$.
For a partition $\lambda\vdash ab$, write $\lambda_2=0$ if it
has one row. Then
\[
 b\geq\lambda_2+\binom a2
 \quad\Longrightarrow\quad
 I(\operatorname{Ch}_b(V))_a[\lambda]=0.                 \tag{A-020.1}
\]
Here $[\lambda]$ denotes the complete $S_\lambda V^*$ isotypic
component, including every multiplicity copy. Consequently the
degree-$d$ coordinate ring of the determinant orbit closure in
$n^2$ variables has full ambient multiplicity at $\lambda$ if
\[
                         n\geq\lambda_2+\binom d2.    \tag{A-020.2}
\]
In this range $K_n(\lambda)=a_\lambda(d[n])$, so no other
orbit closure in the same representation can have a larger
degree-$d$ multiplicity at $\lambda$.

The argument refines the two differential sequences in OpenAI,
\emph{Quadratic stabilization of the canonical Foulkes--Howe map}
(September 25, 2026), Proposition 3, by keeping track of the
diagonal $GL(V)$ type. The uniform bound $b\geq a(a-1)$ from
that source remains a separate, complementary sufficient
condition. No historical-priority claim is made for the present
shape bound.

\paragraph{A differential kernel lemma.}
Let $U,Y$ be independent vector-variable blocks, and set
$D=\sum_iY_i\partial/\partial U_i$. On the polynomial space
of bidegree $(p,q)$, as a diagonal $GL(V)$ module,
\[
 \ker D=
 \begin{cases}
 S_{(q,p)}V^*,&p\leq q,\\
 0,&p>q.
 \end{cases}                                          \tag{A-020.3}
\]
The usual convention makes the first module zero if its
partition has too many rows. If additional independent vector
blocks of fixed nonnegative degrees are present, every
irreducible constituent $S_\eta V^*$ of the kernel satisfies
$\eta_2\geq p$. In particular the operator is injective on
the $\lambda$ isotypic component whenever $\lambda_2<p$,
irrespective of $q$.

\begin{proof}
The horizontal Pieri rule gives the multiplicity-free
decomposition
\[
 \operatorname{Sym}^pV^*\otimes\operatorname{Sym}^qV^*
 =\bigoplus_{i=0}^{\min(p,q)} S_{(p+q-i,i)}V^*.
\]
For dimension at least two, a highest-weight vector for the
$i$-th summand is
\[
 h_i=(U_1Y_2-U_2Y_1)^i U_1^{p-i}Y_1^{q-i}.
\]
The operator kills the displayed determinant and sends $h_i$
to
\[
 (p-i)(U_1Y_2-U_2Y_1)^i
                  U_1^{p-i-1}Y_1^{q-i+1}.
\]
This is nonzero for $i<p$, so equivariance and irreducibility
make $D$ injective on each such summand. The remaining summand
$i=p$, present exactly when $p\leq q$, is killed. This proves
(A-020.3). In dimension one the same conclusion follows from
$D(U^pY^q)=pU^{p-1}Y^{q+1}$; dimension zero is immediate.
The case $p=0$ says that the whole source is killed and agrees
with $S_{(q,0)}=\operatorname{Sym}^q$.

With spectators the kernel is the tensor product of this kernel
and their polynomial space, since tensoring vector spaces is
exact. That spectator space is a tensor product of symmetric
powers of $V^*$. Each successive horizontal Pieri step only
adds boxes to the partition $(q,p)$, so every final second row
has length at least $p$. This proves the claimed restricted
injectivity. The classical rule used here is Macdonald,
Chapter I, (5.16), printed p.~73. \end{proof}

\paragraph{Proof of (A-020.1).}
We induct on $a$. For $a=1$, pure powers span
$\operatorname{Sym}^bV$, so there are no linear equations.
For $a=2$, there are no quadratic equations when $b\geq2$
by the cited Proposition 3; when $b=1$ the Chow cone is all
of $V$. These base cases are stronger than required.

Let $a\geq3$, put $r=a-1$, $m=b-r$, and assume
$b\geq\lambda_2+\binom a2$. Then
\[
 m\geq\lambda_2+\binom r2\geq1.                       \tag{A-020.4}
\]
Suppose $T$ is in the $\lambda$ isotypic component of the
space of symmetric $a$-linear forms on $\operatorname{Sym}^bV$
whose diagonals vanish on products of $b$ vectors. Such forms
are precisely the polarizations of the degree-$a$ equations.
All constructions below are equivariant under the diagonal
$GL(V)$ action, so their images remain in type $\lambda$.

For a product $Q=y_1\cdots y_m$, the binary polynomial
$x^r+ct^r$ splits into $r$ linear factors over $\mathbb C$.
The coefficient of $c$ in the diagonal vanishing identity for
$(x^r+ct^r)Q$ is therefore
\[
                       T(x^rQ,\ldots,x^rQ,t^rQ)=0.     \tag{A-020.5}
\]
Now set $H=T(x^rQ,\ldots,x^rQ,t^b)$, keeping $x,t,y_i$
as vector-variable blocks. Apply
$D_{y_m,t}\cdots D_{y_1,t}$, where
$D_{Y,U}=\sum Y_i\partial_{U_i}$. The result is
$(b!/r!)$ times (A-020.5), hence zero. Before the $i$-th
operator the degree in $t$ is $b-i+1$, while that in the new
$y_i$ is $r$. Their difference is $m-i+1\geq1$.
The $p>q$ case of (A-020.3), with all other blocks as
spectators, makes the composition injective. Thus $H=0$.

Consider the equivariant linear map sending $T$ to the
polynomial family of symmetric $r$-linear forms
\[
 S_{x,t}(Q_1,\ldots,Q_r)
             =T(x^rQ_1,\ldots,x^rQ_r,t^b).
\]
Its codomain is
\[
 \operatorname{Sym}^{r^2}V^*\otimes\operatorname{Sym}^bV^*
             \otimes\operatorname{Sym}^r((\operatorname{Sym}^mV)^*).
\]
The identity $H=0$ says that every coefficient in $x,t$
of its last factor vanishes diagonally on products of $m$
vectors. Consequently the image lies in
\[
 \operatorname{Sym}^{r^2}V^*\otimes\operatorname{Sym}^bV^*
                         \otimes I(\operatorname{Ch}_m(V))_r.
                                                               \tag{A-020.6}
\]
By induction and (A-020.4), every constituent $S_\nu V^*$
of the last factor has
$\nu_2>m-\binom r2=b-\binom a2\geq\lambda_2$.
Multiplication by the two symmetric-power factors in (A-020.6)
is governed by two horizontal Pieri steps. It can only add boxes
to $\nu$, so cannot produce any constituent of type $\lambda$.
Since the map is equivariant and its domain is isotypic of type
$\lambda$, its image is zero. Therefore
\[
 T(x^rQ_1,\ldots,x^rQ_r,t^b)=0
       \quad\hbox{for all }Q_i\in\operatorname{Sym}^mV.       \tag{A-020.7}
\]

Finally introduce independent $v_1,\ldots,v_r,t,x$ and put
$f=T(v_1^b,\ldots,v_r^b,t^b)$. Apply $D_{x,v_i}^r$
successively for $i=1,\ldots,r$. The result is
\[
 \left(\frac{b!}{m!}\right)^r
                 T(x^rv_1^m,\ldots,x^rv_r^m,t^b)=0
\]
by (A-020.7). Before each individual derivative the source
$v_i$ has degree $p\geq b-r+1$. Our hypothesis implies
\[
 \lambda_2\leq b-\binom a2\leq b-r<p.
\]
Every intermediate polynomial remains in type $\lambda$;
the restricted injectivity part of (A-020.3) therefore applies
at every step, even when the common $x$-degree is larger than
the source degree. The composition is injective on that type,
so $f=0$. Pure powers span each input space, hence $T=0$.
This completes the induction and proves (A-020.1).

For (A-020.2), products of $n$ linear forms are determinants
of diagonal matrices of those forms. They are linear
endomorphism substitutions into $\det_n$, and belong to its
orbit closure because invertible endomorphisms are dense.
The determinant-closure ideal is therefore contained in the
Chow ideal. Its $\lambda$ component vanishes by (A-020.1).
Semisimplicity in characteristic zero then shows that every
ambient copy of $S_\lambda V^*$ survives in the quotient
coordinate ring. This is precisely the stated equality of
multiplicities. \hfill$\square$

\paragraph{Grothendieck dictionary and remaining gate.}
The Chow cone can be defined over $\mathbb Z$ as the scheme
closure of the multiplication morphism from $V^b$ to
$\operatorname{Sym}^bV$. Its homogeneous ideal is a
$GL(V)$-comodule; the operators $D_{Y,U}$ are integral,
equivariant differential maps between multigraded polynomial
modules. Representations as comodules are SGA 3, I,
Proposition 4.7.2, pp.~34--35; torus gradings are I,
Proposition 4.7.3, pp.~35--36. The symmetric and Schur
functors live in the exact $K_0$ and special $\lambda$-ring
framework of SGA 6, VI, Theorem 3.3, pp.~392--393.
The isotypic decomposition, nonzero differential coefficients,
and binary factorization in the proof are used over
$\mathbb C$; no assertion of identical kernel ranks in every
characteristic follows from the integral definitions.

Together with M-010, an actual multiplicity obstruction must
satisfy both $d(d-1)>n$ and
$\lambda_2>n-\binom d2$, in addition to the other proved
shape restrictions. The second condition is informative for
degrees with $\binom d2\leq n$, and retains all multiplicity
copies. It concerns the degree of equations and the second row,
not a saturation assertion for Kronecker support.
